% ============================================================
% DISCUSSION
% ============================================================

\section*{Discussion}


% ----------------------------------------------------------
% Classical Thickness Distributions
% ----------------------------------------------------------

\subsection*{Classical Design Approach}

The classical airfoil design approach was primarily theoretical and geometry-driven. 
Engineers defined airfoil shapes using analytical equations without detailed aerodynamic performance predictions. 
The aerodynamic characteristics were then evaluated through experiments and empirical observations, followed by iterative modifications.

Although this approach established important theoretical foundations and provided valuable design concepts, 
it often resulted in inefficient designs and a relatively high failure ratio. The lack of accurate predictive tools made it difficult to optimize airfoils for specific aerodynamic requirements.

% ----------------------------------------------------------
% Classical Thickness Distributions
% ----------------------------------------------------------

\subsection*{Modern Design Approach}

Modern airfoil design relies heavily on numerical methods and computational tools. 
Airfoils are now designed to meet specific aerodynamic performance targets, such as desired lift coefficients $C_l$, lift-to-drag ratios, and pressure or velocity distributions along the chord.

The modern design process typically follows these steps:

\begin{enumerate}
    \item Define aerodynamic performance requirements
    \item Generate airfoil geometry using numerical optimization methods
    \item Simulate aerodynamic performance using computational tools
    \item Manufacture prototypes
    \item Conduct experimental validation
    \item Refine and optimize the design based on results
\end{enumerate}

This approach significantly improves design reliability, efficiency, and performance compared to classical methods.

Modern airfoils are no longer defined solely by geometric parameters but by aerodynamic performance objectives. 
The development of computational fluid dynamics (CFD) and numerical optimization 
allows engineers to design airfoils with precise aerodynamic characteristics.

% ----------------------------------------------------------
% Classical Thickness Distributions
% ----------------------------------------------------------

\subsection*{Limitations of 2D Airfoil Analysis and Transition to 3D Design}

Airfoil sections represent only a two-dimensional approximation of wing behavior. 
The aerodynamic characteristics of a 2D airfoil correspond to those of an infinite wing with an infinite aspect ratio, 
which does not exist in practical applications.

In reality, wings, propeller blades, and turbine blades are three-dimensional and finite structures. 
Their aerodynamic behavior is influenced by additional factors such as wingtip vortices, induced drag, and spanwise flow. 
As a result, accurate aerodynamic analysis requires three-dimensional numerical simulations.

Fortunately, modern aerospace engineering provides advanced computational tools such as OpenFOAM, OpenVSP, and ANSYS, 
which enable engineers to perform detailed aerodynamic simulations and optimize designs efficiently.

The development of such tools reflects the interdisciplinary nature of modern aerospace engineering. 
Their creation requires expertise in aerodynamics, applied mathematics, numerical methods, and computer science. 
This interdisciplinary collaboration has significantly improved the accuracy, efficiency, and reliability of airfoil and aircraft design.